\documentclass[10pt,a4paper]{amsart}
\usepackage{amsmath,amssymb,amsthm,graphicx}
% Retained minimal preamble for the Stacks Project chapter more-algebra.tex (More on Algebra).
% The chapter file pinned-original-source.tex carries only \input{preamble} on its first line, so no
% preamble text lives inside the chapter itself.  The Stacks Project uses one single shared preamble for
% every chapter; the snapshot stages it identically as sources/stacks-chow/preamble.tex and
% sources/stacks-etale-cohomology/preamble.tex.  The block below is copied mechanically, line for line,
% from sources/stacks-chow/preamble.tex lines 14-18 (the Xy-pic setup) and 154-217 (the native theorem
% declarations together with the mathematical macros).  Nothing was added, removed or redefined; the
% wrapper class and the amsmath/amssymb/amsthm/graphicx base packages are supplied by the assembler.
% Chapter: more-algebra.tex -- More on Algebra.
% For commutative diagrams we use Xy-pic
\usepackage[all]{xy}

% We use 2cell for 2-commutative diagrams.
\xyoption{2cell}

% Theorem environments.
%
\theoremstyle{plain}
\newtheorem{theorem}[subsection]{Theorem}
\newtheorem{proposition}[subsection]{Proposition}
\newtheorem{lemma}[subsection]{Lemma}

\theoremstyle{definition}
\newtheorem{definition}[subsection]{Definition}
\newtheorem{example}[subsection]{Example}
\newtheorem{exercise}[subsection]{Exercise}
\newtheorem{situation}[subsection]{Situation}

\theoremstyle{remark}
\newtheorem{remark}[subsection]{Remark}
\newtheorem{remarks}[subsection]{Remarks}

\numberwithin{equation}{subsection}

% Macros
%
\def\lim{\mathop{\mathrm{lim}}\nolimits}
\def\colim{\mathop{\mathrm{colim}}\nolimits}
\def\Spec{\mathop{\mathrm{Spec}}}
\def\Hom{\mathop{\mathrm{Hom}}\nolimits}
\def\Ext{\mathop{\mathrm{Ext}}\nolimits}
\def\SheafHom{\mathop{\mathcal{H}\!\mathit{om}}\nolimits}
\def\SheafExt{\mathop{\mathcal{E}\!\mathit{xt}}\nolimits}
\def\Sch{\mathit{Sch}}
\def\Mor{\mathop{\mathrm{Mor}}\nolimits}
\def\Ob{\mathop{\mathrm{Ob}}\nolimits}
\def\Sh{\mathop{\mathit{Sh}}\nolimits}
\def\NL{\mathop{N\!L}\nolimits}
\def\CH{\mathop{\mathrm{CH}}\nolimits}
\def\proetale{{pro\text{-}\acute{e}tale}}
\def\etale{{\acute{e}tale}}
\def\QCoh{\mathit{QCoh}}
\def\Ker{\mathop{\mathrm{Ker}}}
\def\Im{\mathop{\mathrm{Im}}}
\def\Coker{\mathop{\mathrm{Coker}}}
\def\Coim{\mathop{\mathrm{Coim}}}

% Boxtimes
%
\DeclareMathSymbol{\boxtimes}{\mathbin}{AMSa}{"02}

%
% Macros for moduli stacks/spaces
%
\def\QCohstack{\mathcal{QC}\!\mathit{oh}}
\def\Cohstack{\mathcal{C}\!\mathit{oh}}
\def\Spacesstack{\mathcal{S}\!\mathit{paces}}
\def\Quotfunctor{\mathrm{Quot}}
\def\Hilbfunctor{\mathrm{Hilb}}
\def\Curvesstack{\mathcal{C}\!\mathit{urves}}
\def\Polarizedstack{\mathcal{P}\!\mathit{olarized}}
\def\Complexesstack{\mathcal{C}\!\mathit{omplexes}}
% \Pic is the operator that assigns to X its picard group, usage \Pic(X)
% \Picardstack_{X/B} denotes the Picard stack of X over B
% \Picardfunctor_{X/B} denotes the Picard functor of X over B
\def\Pic{\mathop{\mathrm{Pic}}\nolimits}
\def\Picardstack{\mathcal{P}\!\mathit{ic}}
\def\Picardfunctor{\mathrm{Pic}}
\def\Deformationcategory{\mathcal{D}\!\mathit{ef}}

% --- appended 2026-10-12 (Chief): the project own definitions of the reference/slogan/history
% display environments, copied mechanically from the shared Stacks preamble lines 8-12
% (\usepackage{verbatim} plus the three \newenvironment lines, which the chapter preamble keeps
% verbatim from the project preamble). Nothing added, removed or redefined: these environments
% carry no mathematical content and the project itself typesets them as comments. Without these
% three lines a counted statement whose own text uses them fails with Environment slogan undefined.
\usepackage{verbatim}
\newenvironment{reference}{\comment}{\endcomment}
\newenvironment{slogan}{\comment}{\endcomment}
\newenvironment{history}{\comment}{\endcomment}
\usepackage{cleveref}
\crefname{originallemma}{original lemma}{original lemmas}
\Crefname{originallemma}{Original lemma}{Original lemmas}
\makeatletter
\newlabel{algebra-lemma-characterize-etale}{{00U6}{source}{}{sourcepointer-1}{}}
\newlabel{algebra-lemma-characterize-etale@cref}{{[originallemma][0][]00U6}{[1][0][]source}{}{}{}}
\makeatother
\crefname{originallemma}{original lemma}{original lemmas}
\Crefname{originallemma}{Original lemma}{Original lemmas}
\makeatletter
\newlabel{algebra-lemma-finite-extension-dim-1}{{02MJ}{source}{}{sourcepointer-2}{}}
\newlabel{algebra-lemma-finite-extension-dim-1@cref}{{[originallemma][0][]02MJ}{[1][0][]source}{}{}{}}
\makeatother
\crefname{originallemma}{original lemma}{original lemmas}
\Crefname{originallemma}{Original lemma}{Original lemmas}
\makeatletter
\newlabel{fields-lemma-finite-Galois}{{09I1}{source}{}{sourcepointer-3}{}}
\newlabel{fields-lemma-finite-Galois@cref}{{[originallemma][0][]09I1}{[1][0][]source}{}{}{}}
\makeatother
\crefname{originallemma}{original lemma}{original lemmas}
\Crefname{originallemma}{Original lemma}{Original lemmas}
\makeatletter
\newlabel{fields-lemma-galois-over-fixed-field}{{09I3}{source}{}{sourcepointer-4}{}}
\newlabel{fields-lemma-galois-over-fixed-field@cref}{{[originallemma][0][]09I3}{[1][0][]source}{}{}{}}
\makeatother
\crefname{originallemma}{original lemma}{original lemmas}
\Crefname{originallemma}{Original lemma}{Original lemmas}
\makeatletter
\newlabel{lemma-galois}{{33179}{source}{}{sourcepointer-5}{}}
\newlabel{lemma-galois@cref}{{[originallemma][0][]33179}{[1][0][]source}{}{}{}}
\makeatother
\crefname{originallemma}{original lemma}{original lemmas}
\Crefname{originallemma}{Original lemma}{Original lemmas}
\makeatletter
\newlabel{lemma-invariants-modulo}{{32553}{source}{}{sourcepointer-6}{}}
\newlabel{lemma-invariants-modulo@cref}{{[originallemma][0][]32553}{[1][0][]source}{}{}{}}
\makeatother
\crefname{originallemma}{original lemma}{original lemmas}
\Crefname{originallemma}{Original lemma}{Original lemmas}
\makeatletter
\newlabel{lemma-one-orbit}{{32680}{source}{}{sourcepointer-7}{}}
\newlabel{lemma-one-orbit@cref}{{[originallemma][0][]32680}{[1][0][]source}{}{}{}}
\makeatother
\crefname{originallemma}{original lemma}{original lemmas}
\Crefname{originallemma}{Original lemma}{Original lemmas}
\makeatletter
\newlabel{lemma-one-orbit-geometric}{{32704}{source}{}{sourcepointer-8}{}}
\newlabel{lemma-one-orbit-geometric@cref}{{[originallemma][0][]32704}{[1][0][]source}{}{}{}}
\makeatother
\crefname{originalremark}{original remark}{original remarks}
\Crefname{originalremark}{Original remark}{Original remarks}
\makeatletter
\newlabel{remark-finite-separable-extension}{{33006}{source}{}{sourcepointer-9}{}}
\newlabel{remark-finite-separable-extension@cref}{{[originalremark][0][]33006}{[1][0][]source}{}{}{}}
\makeatother
\title{The invariants of the inertia group form the integral closure of a Galois subextension, and the corresponding local extension is etale}
\author{The Stacks Project Authors}
\date{stacks-project:more-algebra.tex:local-snapshot-20261012; GNU FDL 1.2 or later}
\begin{document}
\maketitle
\setcounter{subsection}{1}
\noindent\textit{Editorial note.} All delivered mathematical text is the author's own, copied from the pinned Stacks Project chapter file more-algebra.tex; no statement or proof step was rewritten, reordered or removed.\par\medskip
\noindent\textit{Editorial note.} The chapter file pinned-original-source.tex carries only \textbackslash{}input\{preamble\} on its first line, so no preamble text lives inside the chapter. The retained file sources/stacks-more-algebra/retained-preamble.tex is a mechanical, line-for-line copy of the project's single shared preamble as staged in the snapshot; nothing was added, removed or redefined and no package was added.\par\medskip
\noindent\textit{Editorial note.} The source URL is the source-level Stacks Project URL, as required when more than one card is produced from one source; the per-claim tag is carried by claim\_key.\par\medskip
\noindent\textit{Editorial note.} The author writes 'M\_1 = B\textasciicircum{}I is the integral closure' inside the counted proof span. The nine references inside the counted proof span are recorded as non-load-bearing navigation pointers; they include cross-chapter tags and same-chapter label lines.\par\medskip
\noindent\textit{Allowed standard prerequisites (editorial).}\begin{itemize}
\item The inertia group of a maximal ideal in a Galois extension of a discretely valued field, and the decomposition group. fixed at statement lines 33477-33479 and used throughout the counted proof span. Supplies the group I and the quotient D/I whose action on the residue field is analysed.
\item The norm of a uniformizer and multiplicativity of ramification indices in a tower of discretely valued fields. used at proof lines 33503-33520. Produces the two degree formulae whose comparison yields the etaleness conclusion.
\item Purely inseparable residue extension for the invariants subgroup, and faithfulness of the action of D/I. used at proof lines 33490-33498. Gives the bound [kappa(m') : kappa] >= |D/I| that the index comparison consumes.
\end{itemize}\medskip
\ifdefined\hypertarget\hypertarget{sourcepointer-1}{}\fi
\noindent\textit{Source reference.} Original lemma 00U6: focused-seeds/AGP/sources/stacks-snapshot-20261011/tags.txt, source lines 1100--1100. Cross-chapter navigation pointer into the Stacks Project; the standard fact it supplies is applied inside the counted proof span and the argument resting on it stays there.\par\medskip
\ifdefined\hypertarget\hypertarget{sourcepointer-2}{}\fi
\noindent\textit{Source reference.} Original lemma 02MJ: focused-seeds/AGP/sources/stacks-snapshot-20261011/tags.txt, source lines 3446--3446. Cross-chapter navigation pointer into the Stacks Project; the standard fact it supplies is applied inside the counted proof span and the argument resting on it stays there.\par\medskip
\ifdefined\hypertarget\hypertarget{sourcepointer-3}{}\fi
\noindent\textit{Source reference.} Original lemma 09I1: focused-seeds/AGP/sources/stacks-snapshot-20261011/tags.txt, source lines 12024--12024. Cross-chapter navigation pointer into the Stacks Project; the standard fact it supplies is applied inside the counted proof span and the argument resting on it stays there.\par\medskip
\ifdefined\hypertarget\hypertarget{sourcepointer-4}{}\fi
\noindent\textit{Source reference.} Original lemma 09I3: focused-seeds/AGP/sources/stacks-snapshot-20261011/tags.txt, source lines 12026--12026. Cross-chapter navigation pointer into the Stacks Project; the standard fact it supplies is applied inside the counted proof span and the argument resting on it stays there.\par\medskip
\ifdefined\hypertarget\hypertarget{sourcepointer-5}{}\fi
\noindent\textit{Source reference.} Original lemma 33179: sources/stacks-more-algebra/pinned-original-source.tex, source lines 33179--33179. Same-chapter navigation pointer; the statement it names is standard for the section and the argument that consumes it stays inside the counted proof span.\par\medskip
\ifdefined\hypertarget\hypertarget{sourcepointer-6}{}\fi
\noindent\textit{Source reference.} Original lemma 32553: sources/stacks-more-algebra/pinned-original-source.tex, source lines 32553--32553. Same-chapter navigation pointer; the statement it names is standard for the section and the argument that consumes it stays inside the counted proof span.\par\medskip
\ifdefined\hypertarget\hypertarget{sourcepointer-7}{}\fi
\noindent\textit{Source reference.} Original lemma 32680: sources/stacks-more-algebra/pinned-original-source.tex, source lines 32680--32680. Same-chapter navigation pointer; the statement it names is standard for the section and the argument that consumes it stays inside the counted proof span.\par\medskip
\ifdefined\hypertarget\hypertarget{sourcepointer-8}{}\fi
\noindent\textit{Source reference.} Original lemma 32704: sources/stacks-more-algebra/pinned-original-source.tex, source lines 32704--32704. Same-chapter navigation pointer; the statement it names is standard for the section and the argument that consumes it stays inside the counted proof span.\par\medskip
\ifdefined\hypertarget\hypertarget{sourcepointer-9}{}\fi
\noindent\textit{Source reference.} Original remark 33006: sources/stacks-more-algebra/pinned-original-source.tex, source lines 33006--33006. Same-chapter navigation pointer; the statement it names is standard for the section and the argument that consumes it stays inside the counted proof span.\par\medskip
\begin{lemma}
\label{lemma-inertial-invariants-unramified}
Let $A$ be a discrete valuation ring with fraction field $K$.
Let $L/K$ be a finite Galois extension. Let $\mathfrak m \subset B$
be a maximal ideal of the integral closure of $A$ in $L$.
Let $I \subset G$ be the inertia group of $\mathfrak m$.
Then $B^I$ is the integral closure of $A$ in $L^I$ and
$A \to (B^I)_{B^I \cap \mathfrak m}$ is \'etale.
\end{lemma}

\begin{proof}
Write $B' = B^I$. It follows from the definitions that $B' = B^I$
is the integral closure of $A$ in $L^I$. Write
$\mathfrak m' = B^I \cap \mathfrak m = B' \cap \mathfrak m \subset B'$.
By Lemma \ref{lemma-one-orbit} the maximal ideal $\mathfrak m$ is the
unique prime ideal of $B$ lying over $\mathfrak m'$.
As $I$ acts trivially on $\kappa(\mathfrak m)$ we see from
Lemma \ref{lemma-invariants-modulo} that the extension
$\kappa(\mathfrak m)/\kappa(\mathfrak m')$ is purely inseparable
(perhaps an easier alternative is to apply the result of
Lemma \ref{lemma-one-orbit-geometric}).
Since $D/I$ acts faithfully on $\kappa(\mathfrak m')$,
we conclude that $D/I$ acts faithfully on $\kappa(\mathfrak m)$.
Of course the elements of the residue field $\kappa$ of $A$
are fixed by this action.
By Galois theory we see that $[\kappa(\mathfrak m') : \kappa] \geq |D/I|$,
see Fields, Lemma \ref{fields-lemma-galois-over-fixed-field}.

\medskip\noindent
Let $\pi$ be the uniformizer of $A$. Since
$\text{Norm}_{L/K}(\pi) = \pi^{[L : K]}$ we see from
Algebra, Lemma \ref{algebra-lemma-finite-extension-dim-1}
that
$$
|G| = [L : K] = [L : K]\ \text{ord}_A(\pi) =
|G/D|\ [\kappa(\mathfrak m) : \kappa]\ \text{ord}_{B_\mathfrak m}(\pi)
$$
as there are $n = |G/D|$ maximal ideals of $B$ which are all
conjugate under $G$, see
Remark \ref{remark-finite-separable-extension} and
Lemma \ref{lemma-galois}.
Applying the same reasoning to the finite extension
the finite extension $L/L^I$ of degree $|I|$ we find
$$
|I|\ \text{ord}_{B'_{\mathfrak m'}}(\pi) =
[\kappa(\mathfrak m) : \kappa(\mathfrak m')]\ \text{ord}_{B_\mathfrak m}(\pi)
$$
We conclude that
$$
\text{ord}_{B'_{\mathfrak m'}}(\pi) =
\frac{|D/I|}{[\kappa(\mathfrak m') : \kappa]}
$$
Since the left hand side is a positive integer and since the right hand
side is $\leq 1$ by the above, we conclude that we have equality,
$\text{ord}_{B'_{\mathfrak m'}}(\pi) = 1$ and
$\kappa(\mathfrak m')/\kappa$ has degree $|D/I|$.
Thus $\pi B'_{\mathfrak m'} = \mathfrak m' B_\mathfrak m'$ and
$\kappa(\mathfrak m')$ is Galois over $\kappa$ with
Galois group $D/I$, in particular separable, see
Fields, Lemma \ref{fields-lemma-finite-Galois}.
By Algebra, Lemma \ref{algebra-lemma-characterize-etale}
we find that $A \to B'_{\mathfrak m'}$ is \'etale
as desired.
\end{proof}

\end{document}
